\subsection{\label{subsec.fps.exp}Exponentials and logarithms}

\begin{convention}
\label{conv.fps.exp.K-Q-alg}Throughout Section \ref{subsec.fps.exp}, we assume
that $K$ is not just a commutative ring, but actually a commutative
$\mathbb{Q}$-algebra.
\end{convention}

\subsubsection{Definitions}

Convention \ref{conv.fps.exp.K-Q-alg} entails that elements of $K$ can be
divided by the positive integers $1,2,3,\ldots$. We can use this to define
three specific (and particularly important) FPSs over $K$:

\begin{definition}
\label{def.fps.exp-log}Define three FPS $\exp$, $\overline{\log}$ and
$\overline{\exp}$ in $K\left[  \left[  x\right]  \right]  $ by%
\begin{align*}
\exp &  :=\sum_{n\in\mathbb{N}}\dfrac{1}{n!}x^{n},\\
\overline{\log}  &  :=\sum_{n\geq1}\dfrac{\left(  -1\right)  ^{n-1}}{n}%
x^{n},\\
\overline{\exp}  &  :=\exp-1=\sum_{n\geq1}\dfrac{1}{n!}x^{n}.
\end{align*}


(The last equality sign here follows from $\exp=\sum_{n\in\mathbb{N}}\dfrac
{1}{n!}x^{n}=\underbrace{\dfrac{1}{0!}}_{=1}\underbrace{x^{0}}_{=1}%
+\sum_{n\geq1}\dfrac{1}{n!}x^{n}=1+\sum_{n\geq1}\dfrac{1}{n!}x^{n}$.)
\end{definition}

Note that the FPS $\exp$ is the usual exponential series from analysis, but
now manifesting itself as a FPS. Likewise, $\overline{\log}$ is the Mercator
series for $\log\left(  1+x\right)  $, where $\log$ stands for the natural
logarithm function. The natural logarithm function itself cannot be
interpreted as an FPS, since $\log0$ is undefined.

\subsubsection{The exponential and the logarithm are inverse}

I will prove that
\begin{equation}
\overline{\exp}\circ\overline{\log}=\overline{\log}\circ\overline{\exp}=x.
\label{eq.fps.exp-log-inverse.1}%
\end{equation}
This is an algebraic analogue of the well-known fact from analysis which
states that the exponential and logarithm functions are mutually inverse.

There is a short way of proving (\ref{eq.fps.exp-log-inverse.1}), which I will
not take: Namely, one can show that any equality between holomorphic functions
on an open disk around the origin leads to an equality between their Taylor
series (viewed as FPSs). Thus, if you have proved in complex analysis that
$\log\circ\exp=\operatorname*{id}$ on an open disk around $0$ and $\exp
\circ\log=\operatorname*{id}$ on an open disk around $1$, then you
automatically get (\ref{eq.fps.exp-log-inverse.1}) (indeed, $\overline{\exp}$
and $\overline{\log}$ are the Taylor series of the functions $\exp$ and $\log$
around $1$, with the reservation that the point $1$ has been moved to the
origin by a shift). This approach uses nontrivial results from complex
analysis, so I will not follow it and instead start from scratch.

The main tool in the proof of (\ref{eq.fps.exp-log-inverse.1}) will be the
following useful proposition (\cite[Lemma 0.4]{logexp}):

\begin{proposition}
\label{prop.fps.exp-log-der}Let $g\in K\left[  \left[  x\right]  \right]  $
with $\left[  x^{0}\right]  g=0$. Then: \medskip

\textbf{(a)} We have%
\[
\left(  \overline{\exp}\circ g\right)  ^{\prime}=\left(  \exp\circ g\right)
^{\prime}=\left(  \exp\circ g\right)  \cdot g^{\prime}.
\]


\textbf{(b)} We have%
\[
\left(  \overline{\log}\circ g\right)  ^{\prime}=\left(  1+g\right)
^{-1}\cdot g^{\prime}.
\]

\end{proposition}

\begin{proof}
[Proof of Proposition \ref{prop.fps.exp-log-der}.]\textbf{(a)} Let us first
show that $\overline{\exp}^{\prime}=\exp^{\prime}=\exp$. Indeed,
$\overline{\exp}=\exp-1$, so that $\exp=\overline{\exp}+1$ and therefore%
\begin{align}
\exp^{\prime}  &  =\left(  \overline{\exp}+1\right)  ^{\prime}=\overline{\exp
}^{\prime}+\underbrace{1^{\prime}}_{=0}\ \ \ \ \ \ \ \ \ \ \left(  \text{by
Theorem \ref{thm.fps.deriv.rules} \textbf{(a)}}\right) \nonumber\\
&  =\overline{\exp}^{\prime}. \label{pf.prop.fps.exp-log-der.a.1}%
\end{align}
Next, we recall that $\exp=\sum_{n\in\mathbb{N}}\dfrac{1}{n!}x^{n}$. Hence,
the definition of a derivative yields%
\begin{align}
\exp^{\prime}  &  =\sum_{n\geq1}\underbrace{n\cdot\dfrac{1}{n!}}%
_{\substack{=\dfrac{1}{\left(  n-1\right)  !}\\\text{(since }n!=n\cdot\left(
n-1\right)  !\text{)}}}x^{n-1}=\sum_{n\geq1}\dfrac{1}{\left(  n-1\right)
!}x^{n-1}=\sum_{n\in\mathbb{N}}\dfrac{1}{n!}x^{n}\nonumber\\
&  \ \ \ \ \ \ \ \ \ \ \ \ \ \ \ \ \ \ \ \ \left(  \text{here, we have
substituted }n\text{ for }n-1\text{ in the sum}\right) \nonumber\\
&  =\exp. \label{pf.prop.fps.exp-log-der.a.2}%
\end{align}
Comparing this with (\ref{pf.prop.fps.exp-log-der.a.1}), we find
\begin{equation}
\overline{\exp}^{\prime}=\exp. \label{pf.prop.fps.exp-log-der.a.3}%
\end{equation}


Now, we can apply the chain rule (Theorem \ref{thm.fps.deriv.rules}
\textbf{(g)}) to $f=\overline{\exp}$ (since $\left[  x^{0}\right]  g=0$), and
thus obtain%
\[
\left(  \overline{\exp}\circ g\right)  ^{\prime}=\left(  \underbrace{\overline
{\exp}^{\prime}}_{=\exp}\circ\,g\right)  \cdot g^{\prime}=\left(  \exp\circ
g\right)  \cdot g^{\prime}.
\]
The same computation (but with $\overline{\exp}$ replaced by $\exp$) yields
$\left(  \exp\circ g\right)  ^{\prime}=\left(  \exp\circ g\right)  \cdot
g^{\prime}$. Combining these two formulas, we obtain $\left(  \overline{\exp
}\circ g\right)  ^{\prime}=\left(  \exp\circ g\right)  ^{\prime}=\left(
\exp\circ g\right)  \cdot g^{\prime}$. Thus, we have proved Proposition
\ref{prop.fps.exp-log-der} \textbf{(a)}.

\textbf{(b)} We have $\overline{\log}=\sum_{n\geq1}\dfrac{\left(  -1\right)
^{n-1}}{n}x^{n}$. Thus,
\begin{align*}
\overline{\log}^{\prime}  &  =\left(  \sum_{n\geq1}\dfrac{\left(  -1\right)
^{n-1}}{n}x^{n}\right)  ^{\prime}\\
&  =\sum_{n\geq1}\dfrac{\left(  -1\right)  ^{n-1}}{n}\underbrace{\left(
x^{n}\right)  ^{\prime}}_{\substack{=nx^{\prime}x^{n-1}\\\text{(by Theorem
\ref{thm.fps.deriv.rules} \textbf{(f)},}\\\text{applied to }x\text{ instead of
}g\text{)}}}\ \ \ \ \ \ \ \ \ \ \left(  \text{by Theorem
\ref{thm.fps.deriv.rules} \textbf{(b)}}\right) \\
&  =\sum_{n\geq1}\underbrace{\dfrac{\left(  -1\right)  ^{n-1}}{n}n}_{=\left(
-1\right)  ^{n-1}}\underbrace{x^{\prime}}_{=1}x^{n-1}=\sum_{n\geq1}\left(
-1\right)  ^{n-1}x^{n-1}=\sum_{n\in\mathbb{N}}\left(  -1\right)  ^{n}x^{n}%
\end{align*}
(here, we have substituted $n$ for $n-1$ in the sum). On the other hand,
Proposition \ref{prop.fps.invertible} yields $\left(  1+x\right)  ^{-1}%
=\sum_{n\in\mathbb{N}}\left(  -1\right)  ^{n}x^{n}$. Comparing these two
equalities, we find%
\begin{equation}
\overline{\log}^{\prime}=\left(  1+x\right)  ^{-1}.
\label{pf.prop.fps.exp-log-der.b.2}%
\end{equation}


Now, we can apply the chain rule (Theorem \ref{thm.fps.deriv.rules}
\textbf{(g)}) to $f=\overline{\log}$ (since $\left[  x^{0}\right]  g=0$), and
thus obtain%
\begin{equation}
\left(  \overline{\log}\circ g\right)  ^{\prime}=\left(  \underbrace{\overline
{\log}^{\prime}}_{=\left(  1+x\right)  ^{-1}}\circ\,g\right)  \cdot g^{\prime
}=\left(  \left(  1+x\right)  ^{-1}\circ g\right)  \cdot g^{\prime}.
\label{pf.prop.fps.exp-log-der.b.3}%
\end{equation}


However, we claim that $\left(  1+x\right)  ^{-1}\circ g=\left(  1+g\right)
^{-1}$. Indeed, Proposition \ref{prop.fps.subs.rules} \textbf{(c)} (applied to
$f_{1}=1$ and $f_{2}=1+x$) yields $\dfrac{1}{1+x}\circ g=\dfrac{1\circ
g}{\left(  1+x\right)  \circ g}$ (since the FPS $1+x$ is invertible). In view
of $\dfrac{1}{1+x}=\left(  1+x\right)  ^{-1}$ and
\[
\underbrace{1}_{=\underline{1}}\circ\,g=\underline{1}\circ g=\underline{1}%
\ \ \ \ \ \ \ \ \ \ \left(  \text{by Proposition \ref{prop.fps.subs.rules}
\textbf{(f)}, applied to }a=1\right)
\]
and%
\[
\left(  1+x\right)  \circ g=1+g\ \ \ \ \ \ \ \ \ \ \left(  \text{this follows
easily from Definition \ref{def.fps.subs}}\right)  ,
\]
this rewrites as $\left(  1+x\right)  ^{-1}\circ g=\dfrac{\underline{1}}%
{1+g}=\left(  1+g\right)  ^{-1}$. Hence, (\ref{pf.prop.fps.exp-log-der.b.3})
becomes%
\[
\left(  \overline{\log}\circ g\right)  ^{\prime}=\underbrace{\left(  \left(
1+x\right)  ^{-1}\circ g\right)  }_{=\left(  1+g\right)  ^{-1}}\cdot
\,g^{\prime}=\left(  1+g\right)  ^{-1}\cdot g^{\prime}.
\]
Thus, we have proved Proposition \ref{prop.fps.exp-log-der} \textbf{(b)}.
\end{proof}

We will need a very simple lemma, which says (in particular) that if two FPSs
have constant terms $0$, then so does their composition:

\begin{lemma}
\label{lem.fps.compos-cst-term-0}Let $f,g\in K\left[  \left[  x\right]
\right]  $ be two FPSs with $\left[  x^{0}\right]  g=0$. Then, $\left[
x^{0}\right]  \left(  f\circ g\right)  =\left[  x^{0}\right]  f$.
\end{lemma}

\begin{proof}
[Proof of Lemma \ref{lem.fps.compos-cst-term-0}.]Write $f$ in the form
$f=\sum_{n\in\mathbb{N}}f_{n}x^{n}$ with $f_{0},f_{1},f_{2},\ldots\in K$.
Thus, $f_{0}=\left[  x^{0}\right]  f$. Now, Definition \ref{def.fps.subs}
yields $f\left[  g\right]  =\sum_{n\in\mathbb{N}}f_{n}g^{n}$. However,
Proposition \ref{prop.fps.subs.wd} \textbf{(c)} yields $\left[  x^{0}\right]
\left(  \sum_{n\in\mathbb{N}}f_{n}g^{n}\right)  =f_{0}=\left[  x^{0}\right]
f$. In view of $f\circ g=f\left[  g\right]  =\sum_{n\in\mathbb{N}}f_{n}g^{n}$,
we can rewrite this as $\left[  x^{0}\right]  \left(  f\circ g\right)
=\left[  x^{0}\right]  f$. This proves Lemma \ref{lem.fps.compos-cst-term-0}.
\end{proof}

Now, we can prove the equalities we promised (\cite[Theorem 0.1]{logexp}):

\begin{theorem}
\label{thm.fps.exp-log-inv}We have%
\[
\overline{\exp}\circ\overline{\log}=x\ \ \ \ \ \ \ \ \ \ \text{and}%
\ \ \ \ \ \ \ \ \ \ \overline{\log}\circ\overline{\exp}=x.
\]

\end{theorem}

\begin{proof}
[Proof of Theorem \ref{thm.fps.exp-log-inv}.]Let us first prove that
$\overline{\log}\circ\overline{\exp}=x$.

Indeed, the idea of this proof is to show that $\overline{\log}\circ
\overline{\exp}$ and $x$ are two FPSs with the same constant term (namely,
$0$) and with the same derivative. Once this is proved, Theorem
\ref{thm.fps.deriv.rules} \textbf{(h)} will easily yield that they are equal.

Let us fill in the details. We have $\left[  x^{0}\right]  \overline{\exp}=0$
(since $\overline{\exp}=\sum_{n\geq1}\dfrac{1}{n!}x^{n}$). Hence, Lemma
\ref{lem.fps.compos-cst-term-0} (applied to $f=\overline{\log}$ and
$g=\overline{\exp}$) yields $\left[  x^{0}\right]  \left(  \overline{\log
}\circ\overline{\exp}\right)  =\left[  x^{0}\right]  \overline{\log}=0$ (since
$\overline{\log}=\sum_{n\geq1}\dfrac{\left(  -1\right)  ^{n-1}}{n}x^{n}$).
Now, (\ref{pf.thm.fps.ring.xn(a-b)=}) yields%
\begin{equation}
\left[  x^{0}\right]  \left(  \overline{\log}\circ\overline{\exp}-x\right)
=\underbrace{\left[  x^{0}\right]  \left(  \overline{\log}\circ\overline{\exp
}\right)  }_{=0}-\underbrace{\left[  x^{0}\right]  x}_{=0}=0.
\label{pf.thm.fps.exp-log-inv.a.1}%
\end{equation}
However, $\overline{\exp}=\exp-1$ and thus $1+\overline{\exp}=\exp$. Now,
Proposition \ref{prop.fps.exp-log-der} \textbf{(b)} (applied to $g=\overline
{\exp}$) yields
\[
\left(  \overline{\log}\circ\overline{\exp}\right)  ^{\prime}=\left(
\underbrace{1+\overline{\exp}}_{=\exp}\right)  ^{-1}\cdot\underbrace{\overline
{\exp}^{\prime}}_{\substack{=\exp\\\text{(by
(\ref{pf.prop.fps.exp-log-der.a.3}))}}}=\exp^{-1}\cdot\exp=1=x^{\prime}%
\]
(since $x^{\prime}=1$). Hence, Theorem \ref{thm.fps.deriv.rules} \textbf{(h)}
(applied to $f=\overline{\log}\circ\overline{\exp}$ and $g=x$) yields that
$\overline{\log}\circ\overline{\exp}-x$ is constant. In other words,
$\overline{\log}\circ\overline{\exp}-x=\underline{a}$ for some $a\in K$.
Consider this $a$. From $\overline{\log}\circ\overline{\exp}-x=\underline{a}$,
we obtain $\left[  x^{0}\right]  \left(  \overline{\log}\circ\overline{\exp
}-x\right)  =\left[  x^{0}\right]  \underline{a}=a$. Comparing this with
(\ref{pf.thm.fps.exp-log-inv.a.1}), we find $a=0$. Hence, $\overline{\log
}\circ\overline{\exp}-x=\underline{a}$ rewrites as $\overline{\log}%
\circ\overline{\exp}-x=\underline{0}$. In other words, $\overline{\log}%
\circ\overline{\exp}=x$.

Now it remains to prove that $\overline{\exp}\circ\overline{\log}=x$. There
are (at least) two ways to do this:

\begin{itemize}
\item \textit{1st way:} A homework exercise (Exercise \ref{exe.fps.comp-inv})
says that any FPS $f$ with $\left[  x^{0}\right]  f=0$ and with $\left[
x^{1}\right]  f$ invertible has a unique compositional inverse (i.e., there is
a unique FPS $g$ with $\left[  x^{0}\right]  g=0$ and $f\circ g=g\circ f=x$).
We can apply this to $f=\overline{\log}$ (since $\left[  x^{0}\right]
\overline{\log}=0$ and since $\left[  x^{1}\right]  \overline{\log}=1$ is
invertible), and thus see that $\overline{\log}$ has a unique compositional
inverse $g$. This compositional inverse $g$ must be $\overline{\exp}$, since
$\overline{\log}\circ\overline{\exp}=x$ (indeed, comparing
$\underbrace{\left(  g\circ\overline{\log}\right)  }_{=x}\circ\,\overline
{\exp}=x\circ\overline{\exp}=\overline{\exp}$ with
\begin{align*}
\left(  g\circ\overline{\log}\right)  \circ\overline{\exp}  &  =g\circ
\underbrace{\left(  \overline{\log}\circ\overline{\exp}\right)  }%
_{=x}\ \ \ \ \ \ \ \ \ \ \left(  \text{by Proposition
\ref{prop.fps.subs.rules} \textbf{(e)}}\right) \\
&  =g\circ x=g
\end{align*}
yields $g=\overline{\exp}$). But this entails that $\overline{\exp}%
\circ\overline{\log}=x$ as well.

\item \textit{2nd way:} Here is a more direct argument. We shall first show
that $\exp\circ\overline{\log}=1+x$.

To wit: The FPS $1+x$ is invertible (by Proposition
\ref{prop.fps.invertible.1+x}). Thus, applying the quotient rule (Theorem
\ref{thm.fps.deriv.rules} \textbf{(e)}) to $f=\exp\circ\,\overline{\log}$ and
$g=1+x$, we obtain%
\[
\left(  \dfrac{\exp\circ\,\overline{\log}}{1+x}\right)  ^{\prime}%
=\dfrac{\left(  \exp\circ\,\overline{\log}\right)  ^{\prime}\cdot\left(
1+x\right)  -\left(  \exp\circ\,\overline{\log}\right)  \cdot\left(
1+x\right)  ^{\prime}}{\left(  1+x\right)  ^{2}}.
\]
In view of%
\begin{align*}
\left(  \exp\circ\,\overline{\log}\right)  ^{\prime}  &  =\left(  \exp
\circ\,\overline{\log}\right)  \cdot\underbrace{\overline{\log}^{\prime}%
}_{\substack{=\left(  1+x\right)  ^{-1}\\\text{(by
(\ref{pf.prop.fps.exp-log-der.b.2}))}}}\\
&  \ \ \ \ \ \ \ \ \ \ \ \ \ \ \ \ \ \ \ \ \left(  \text{by Proposition
\ref{prop.fps.exp-log-der} \textbf{(a)}, applied to }g=\overline{\log}\right)
\\
&  =\left(  \exp\circ\,\overline{\log}\right)  \cdot\left(  1+x\right)  ^{-1}%
\end{align*}
and $\left(  1+x\right)  ^{\prime}=1$, we can rewrite this as%
\begin{align*}
\left(  \dfrac{\exp\circ\,\overline{\log}}{1+x}\right)  ^{\prime}  &
=\dfrac{\left(  \exp\circ\,\overline{\log}\right)  \cdot\left(  1+x\right)
^{-1}\cdot\left(  1+x\right)  -\left(  \exp\circ\,\overline{\log}\right)
\cdot1}{\left(  1+x\right)  ^{2}}\\
&  =\dfrac{\left(  \exp\circ\,\overline{\log}\right)  -\left(  \exp
\circ\,\overline{\log}\right)  }{\left(  1+x\right)  ^{2}}=0=0^{\prime}.
\end{align*}
Thus, Theorem \ref{thm.fps.deriv.rules} \textbf{(h)} (applied to
$f=\dfrac{\exp\circ\,\overline{\log}}{1+x}$ and $g=0$) yields that
$\dfrac{\exp\circ\,\overline{\log}}{1+x}-0$ is constant. In other words,
$\dfrac{\exp\circ\,\overline{\log}}{1+x}$ is constant. In other words,
$\dfrac{\exp\circ\,\overline{\log}}{1+x}=\underline{a}$ for some $a\in K$.
Consider this $a$. From $\dfrac{\exp\circ\,\overline{\log}}{1+x}%
=\underline{a}$, we obtain $\exp\circ\,\overline{\log}=\underline{a}\left(
1+x\right)  =a\left(  1+x\right)  $. Thus,%
\[
\left[  x^{0}\right]  \left(  \exp\circ\,\overline{\log}\right)  =\left[
x^{0}\right]  \left(  a\left(  1+x\right)  \right)  =a.
\]
However, it is easy to see that $\left[  x^{0}\right]  \left(  \exp
\circ\,\overline{\log}\right)  =1$\ \ \ \ \footnote{\textit{Proof.} Recall
that $\left[  x^{0}\right]  \overline{\log}=0$. Hence, Lemma
\ref{lem.fps.compos-cst-term-0} (applied to $f=\exp$ and $g=\overline{\log}$)
yields
\begin{align*}
\left[  x^{0}\right]  \left(  \exp\circ\,\overline{\log}\right)   &  =\left[
x^{0}\right]  \exp=\dfrac{1}{0!}\ \ \ \ \ \ \ \ \ \ \left(  \text{since }%
\exp=\sum_{n\in\mathbb{N}}\dfrac{1}{n!}x^{n}\right) \\
&  =\dfrac{1}{1}=1.
\end{align*}
}. Comparing these two equalities, we find $a=1$. Thus, $\exp\circ
\,\overline{\log}=\underbrace{a}_{=1}\left(  1+x\right)  =1+x$.

Now, $\overline{\exp}=\exp-1=\exp+\,\underline{-1}$. Hence,
\begin{align*}
&  \overline{\exp}\circ\overline{\log}\\
&  =\left(  \exp+\,\underline{-1}\right)  \circ\overline{\log}%
=\underbrace{\exp\circ\,\overline{\log}}_{=1+x}+\underbrace{\underline{-1}%
\circ\overline{\log}}_{\substack{=\underline{-1}\\\text{(by Proposition
\ref{prop.fps.subs.rules} \textbf{(f)},}\\\text{applied to }-1\text{ and
}\overline{\log}\text{ instead of }a\text{ and }g\text{)}}}\\
&  \ \ \ \ \ \ \ \ \ \ \ \ \ \ \ \ \ \ \ \ \left(
\begin{array}
[c]{c}%
\text{by Proposition \ref{prop.fps.subs.rules} \textbf{(a)},}\\
\text{applied to }f_{1}=\exp\text{ and }f_{2}=\underline{-1}\text{ and
}g=\overline{\log}%
\end{array}
\right) \\
&  =1+x+\underline{-1}=\underbrace{\left(  1+\underline{-1}\right)  }%
_{=0}+\,x=x.
\end{align*}

\end{itemize}

Either way, we have shown that $\overline{\exp}\circ\overline{\log}=x$. Thus,
the proof of Theorem \ref{thm.fps.exp-log-inv} is complete.
\end{proof}

\subsubsection{The exponential and the logarithm of an FPS}

In Definition \ref{def.fps.exp-log}, we have found algebraic versions of the
exponential and logarithm functions as FPSs. Next, we shall define analogues
of these functions as operators acting on FPSs (i.e., analogues not of the
functions $\exp$ and $\log$ themselves, but rather of composition with these functions):

\begin{definition}
\label{def.fps.Exp-Log-maps}\textbf{(a)} We let $K\left[  \left[  x\right]
\right]  _{0}$ denote the set of all FPSs $f\in K\left[  \left[  x\right]
\right]  $ with $\left[  x^{0}\right]  f=0$. \medskip

\textbf{(b)} We let $K\left[  \left[  x\right]  \right]  _{1}$ denote the set
of all FPSs $f\in K\left[  \left[  x\right]  \right]  $ with $\left[
x^{0}\right]  f=1$. \medskip

\textbf{(c)} We define two maps%
\begin{align*}
\operatorname*{Exp}:K\left[  \left[  x\right]  \right]  _{0}  &  \rightarrow
K\left[  \left[  x\right]  \right]  _{1},\\
g  &  \mapsto\exp\circ g
\end{align*}
and%
\begin{align*}
\operatorname*{Log}:K\left[  \left[  x\right]  \right]  _{1}  &  \rightarrow
K\left[  \left[  x\right]  \right]  _{0},\\
f  &  \mapsto\overline{\log}\circ\left(  f-1\right)  .
\end{align*}
(These two maps are well-defined according to parts \textbf{(c)} and
\textbf{(d)} of Lemma \ref{lem.fps.Exp-Log-maps-wd} below.)
\end{definition}

The maps $\operatorname*{Exp}$ and $\operatorname*{Log}$ are algebraic
analogues of the maps in complex analysis that take any holomorphic function
$f$ to its exponential and logarithm, respectively (at least within certain
regions in which these things are well-defined). As one would hope, and as we
will soon see, they are mutually inverse. Let us first check that their
definition is justified:

\begin{lemma}
\label{lem.fps.Exp-Log-maps-wd}\textbf{(a)} For any $f,g\in K\left[  \left[
x\right]  \right]  _{0}$, we have $f\circ g\in K\left[  \left[  x\right]
\right]  _{0}$. \medskip

\textbf{(b)} For any $f\in K\left[  \left[  x\right]  \right]  _{1}$ and $g\in
K\left[  \left[  x\right]  \right]  _{0}$, we have $f\circ g\in K\left[
\left[  x\right]  \right]  _{1}$. \medskip

\textbf{(c)} For any $g\in K\left[  \left[  x\right]  \right]  _{0}$, we have
$\exp\circ g\in K\left[  \left[  x\right]  \right]  _{1}$. \medskip

\textbf{(d)} For any $f\in K\left[  \left[  x\right]  \right]  _{1}$, we have
$f-1\in K\left[  \left[  x\right]  \right]  _{0}$ and $\overline{\log}%
\circ\left(  f-1\right)  \in K\left[  \left[  x\right]  \right]  _{0}$.
\end{lemma}

\begin{proof}
[Proof of Lemma \ref{lem.fps.Exp-Log-maps-wd}.]\textbf{(a)} Let $f,g\in
K\left[  \left[  x\right]  \right]  _{0}$. In view of the definition of
$K\left[  \left[  x\right]  \right]  _{0}$, this entails that $\left[
x^{0}\right]  f=0$ and $\left[  x^{0}\right]  g=0$. Hence, Lemma
\ref{lem.fps.compos-cst-term-0} yields $\left[  x^{0}\right]  \left(  f\circ
g\right)  =\left[  x^{0}\right]  f=0$. In other words, $f\circ g\in K\left[
\left[  x\right]  \right]  _{0}$ (by the definition of $K\left[  \left[
x\right]  \right]  _{0}$). This proves Lemma \ref{lem.fps.Exp-Log-maps-wd}
\textbf{(a)}.

\textbf{(b)} This is analogous to the proof of Lemma
\ref{lem.fps.Exp-Log-maps-wd} \textbf{(a)}.

\textbf{(c)} Let $g\in K\left[  \left[  x\right]  \right]  _{0}$. From
$\exp=\sum_{n\in\mathbb{N}}\dfrac{1}{n!}x^{n}$, we obtain $\left[
x^{0}\right]  \exp=\dfrac{1}{0!}=1$, so that $\exp\in K\left[  \left[
x\right]  \right]  _{1}$. Hence, Lemma \ref{lem.fps.Exp-Log-maps-wd}
\textbf{(b)} (applied to $f=\exp$) yields $\exp\circ g\in K\left[  \left[
x\right]  \right]  _{1}$. This proves Lemma \ref{lem.fps.Exp-Log-maps-wd}
\textbf{(c)}.

\textbf{(d)} Let $f\in K\left[  \left[  x\right]  \right]  _{1}$. Thus,
$\left[  x^{0}\right]  f=1$. Now, (\ref{pf.thm.fps.ring.xn(a-b)=}) yields
$\left[  x^{0}\right]  \left(  f-1\right)  =\underbrace{\left[  x^{0}\right]
f}_{=1}-\underbrace{\left[  x^{0}\right]  1}_{=1}=1-1=0$, so that $f-1\in
K\left[  \left[  x\right]  \right]  _{0}$. Furthermore, $\left[  x^{0}\right]
\overline{\log}=0$ (since $\overline{\log}=\sum_{n\geq1}\dfrac{\left(
-1\right)  ^{n-1}}{n}x^{n}$) and thus $\overline{\log}\in K\left[  \left[
x\right]  \right]  _{0}$. Hence, Lemma \ref{lem.fps.Exp-Log-maps-wd}
\textbf{(a)} (applied to $\overline{\log}$ and $f-1$ instead of $f$ and $g$)
yields $\overline{\log}\circ\left(  f-1\right)  \in K\left[  \left[  x\right]
\right]  _{0}$. Thus, Lemma \ref{lem.fps.Exp-Log-maps-wd} \textbf{(d)} is proven.
\end{proof}

\begin{lemma}
\label{lem.fps.Exp-Log-maps-inv}The maps $\operatorname*{Exp}$ and
$\operatorname*{Log}$ are mutually inverse bijections between $K\left[
\left[  x\right]  \right]  _{0}$ and $K\left[  \left[  x\right]  \right]
_{1}$.
\end{lemma}

\begin{proof}
[Proof of Lemma \ref{lem.fps.Exp-Log-maps-inv}.]First, we make a simple
auxiliary observation: Each $g\in K\left[  \left[  x\right]  \right]  _{0}$
satisfies\footnote{As before, the \textquotedblleft$\circ$\textquotedblright%
\ operation behaves like multiplication in the sense of PEMDAS conventions.
Thus, the expression \textquotedblleft$\overline{\exp}\circ g+1$%
\textquotedblright\ means $\left(  \overline{\exp}\circ g\right)  +1$.}%
\begin{equation}
\exp\circ g=\overline{\exp}\circ g+1. \label{pf.lem.fps.Exp-Log-maps-inv.1}%
\end{equation}


[\textit{Proof of (\ref{pf.lem.fps.Exp-Log-maps-inv.1}):} Let $g\in K\left[
\left[  x\right]  \right]  _{0}$. Recall that $\overline{\exp}=\exp-1$, so
that $\exp=\overline{\exp}+1=\overline{\exp}+\underline{1}$. Hence,%
\[
\exp\circ g=\left(  \overline{\exp}+\underline{1}\right)  \circ g=\overline
{\exp}\circ g+\underline{1}\circ g
\]
(by Proposition \ref{prop.fps.subs.rules} \textbf{(a)}, applied to
$f_{1}=\overline{\exp}$ and $f_{2}=\underline{1}$). However, Proposition
\ref{prop.fps.subs.rules} \textbf{(f)} (applied to $a=1$) yields
$\underline{1}\circ g=\underline{1}=1$. Hence, $\exp\circ g=\overline{\exp
}\circ g+\underbrace{\underline{1}\circ g}_{=1}=\overline{\exp}\circ g+1$.
This proves (\ref{pf.lem.fps.Exp-Log-maps-inv.1}).]

Now, let us show that $\operatorname*{Exp}\circ\operatorname*{Log}%
=\operatorname*{id}$. Indeed, we fix some $f\in K\left[  \left[  x\right]
\right]  _{1}$. Then, $f-1\in K\left[  \left[  x\right]  \right]  _{0}$ (by
Lemma \ref{lem.fps.Exp-Log-maps-wd} \textbf{(d)}). Hence, Proposition
\ref{prop.fps.subs.rules} \textbf{(e)} (applied to $\overline{\exp}$,
$\overline{\log}$ and $f-1$ instead of $f$, $g$ and $h$) yields $\left(
\overline{\exp}\circ\overline{\log}\right)  \circ\left(  f-1\right)
=\overline{\exp}\circ\left(  \overline{\log}\circ\left(  f-1\right)  \right)
$. Thus,%
\begin{align}
\overline{\exp}\circ\left(  \overline{\log}\circ\left(  f-1\right)  \right)
&  =\underbrace{\left(  \overline{\exp}\circ\overline{\log}\right)
}_{\substack{=x\\\text{(by Theorem \ref{thm.fps.exp-log-inv})}}}\circ\left(
f-1\right)  =x\circ\left(  f-1\right) \nonumber\\
&  =f-1 \label{pf.lem.fps.Exp-Log-maps-inv.a.1}%
\end{align}
(by Proposition \ref{prop.fps.subs.rules} \textbf{(g)}, applied to $g=f-1$).
However,%
\begin{align*}
\left(  \operatorname*{Exp}\circ\operatorname*{Log}\right)  \left(  f\right)
&  =\operatorname*{Exp}\left(  \operatorname*{Log}f\right) \\
&  =\exp\circ\left(  \operatorname*{Log}f\right)  \ \ \ \ \ \ \ \ \ \ \left(
\text{by the definition of }\operatorname*{Exp}\right) \\
&  =\overline{\exp}\circ\underbrace{\left(  \operatorname*{Log}f\right)
}_{\substack{=\overline{\log}\circ\left(  f-1\right)  \\\text{(by the
definition of }\operatorname*{Log}\text{)}}}+\,1\\
&  \ \ \ \ \ \ \ \ \ \ \ \ \ \ \ \ \ \ \ \ \left(  \text{by
(\ref{pf.lem.fps.Exp-Log-maps-inv.1}), applied to }g=\operatorname*{Log}%
f\right) \\
&  =\underbrace{\overline{\exp}\circ\left(  \overline{\log}\circ\left(
f-1\right)  \right)  }_{\substack{=f-1\\\text{(by
(\ref{pf.lem.fps.Exp-Log-maps-inv.a.1}))}}}+\,1\\
&  =\left(  f-1\right)  +1=f=\operatorname*{id}\left(  f\right)  .
\end{align*}


Forget that we fixed $f$. We thus have shown that $\left(  \operatorname*{Exp}%
\circ\operatorname*{Log}\right)  \left(  f\right)  =\operatorname*{id}\left(
f\right)  $ for each $f\in K\left[  \left[  x\right]  \right]  _{1}$. In other
words, $\operatorname*{Exp}\circ\operatorname*{Log}=\operatorname*{id}$.

Using a similar argument, we can show that $\operatorname*{Log}\circ
\operatorname*{Exp}=\operatorname*{id}$. Indeed, let us fix some $g\in
K\left[  \left[  x\right]  \right]  _{0}$. Hence, Proposition
\ref{prop.fps.subs.rules} \textbf{(e)} (applied to $\overline{\log}$,
$\overline{\exp}$ and $g$ instead of $f$, $g$ and $h$) yields $\left(
\overline{\log}\circ\overline{\exp}\right)  \circ g=\overline{\log}%
\circ\left(  \overline{\exp}\circ g\right)  $. Thus,%
\begin{align}
\overline{\log}\circ\left(  \overline{\exp}\circ g\right)   &
=\underbrace{\left(  \overline{\log}\circ\overline{\exp}\right)
}_{\substack{=x\\\text{(by Theorem \ref{thm.fps.exp-log-inv})}}}\circ
\,g=x\circ g\nonumber\\
&  =g \label{pf.lem.fps.Exp-Log-maps-inv.b.1}%
\end{align}
(by Proposition \ref{prop.fps.subs.rules} \textbf{(g)}). But the definition of
$\operatorname*{Exp}$ yields $\operatorname*{Exp}g=\exp\circ g=\overline{\exp
}\circ g+1$ (by (\ref{pf.lem.fps.Exp-Log-maps-inv.1})). Hence,
$\operatorname*{Exp}g-1=\overline{\exp}\circ g$. Now,%
\begin{align*}
\left(  \operatorname*{Log}\circ\operatorname*{Exp}\right)  \left(  g\right)
&  =\operatorname*{Log}\left(  \operatorname*{Exp}g\right) \\
&  =\overline{\log}\circ\left(  \operatorname*{Exp}g-1\right)
\ \ \ \ \ \ \ \ \ \ \left(  \text{by the definition of }\operatorname*{Log}%
\right) \\
&  =\overline{\log}\circ\left(  \overline{\exp}\circ g\right)
\ \ \ \ \ \ \ \ \ \ \left(  \text{since }\operatorname*{Exp}g-1=\overline
{\exp}\circ g\right) \\
&  =g\ \ \ \ \ \ \ \ \ \ \left(  \text{by
(\ref{pf.lem.fps.Exp-Log-maps-inv.b.1})}\right) \\
&  =\operatorname*{id}\left(  g\right)  .
\end{align*}


Forget that we fixed $g$. We thus have shown that $\left(  \operatorname*{Log}%
\circ\operatorname*{Exp}\right)  \left(  g\right)  =\operatorname*{id}\left(
g\right)  $ for each $g\in K\left[  \left[  x\right]  \right]  _{0}$. In other
words, $\operatorname*{Log}\circ\operatorname*{Exp}=\operatorname*{id}$.
Combining this with $\operatorname*{Exp}\circ\operatorname*{Log}%
=\operatorname*{id}$, we see that the maps $\operatorname*{Exp}$ and
$\operatorname*{Log}$ are mutually inverse bijections between $K\left[
\left[  x\right]  \right]  _{0}$ and $K\left[  \left[  x\right]  \right]
_{1}$. This proves Lemma \ref{lem.fps.Exp-Log-maps-inv}.
\end{proof}

\subsubsection{Addition to multiplication}

We will now prove another lemma, which says that the $\operatorname*{Exp}$ and
$\operatorname*{Log}$ maps deserve their names:

\begin{lemma}
\label{lem.fps.Exp-Log-additive}\textbf{(a)} For any $f,g\in K\left[  \left[
x\right]  \right]  _{0}$, we have%
\[
\operatorname*{Exp}\left(  f+g\right)  =\operatorname*{Exp}f\cdot
\operatorname*{Exp}g.
\]


\textbf{(b)} For any $f,g\in K\left[  \left[  x\right]  \right]  _{1}$, we
have%
\[
\operatorname*{Log}\left(  fg\right)  =\operatorname*{Log}%
f+\operatorname*{Log}g.
\]

\end{lemma}

\begin{proof}
[Proof of Lemma \ref{lem.fps.Exp-Log-additive} (sketched).]\textbf{(a)} Like
many of our arguments involving FPSs, this will be a short computation
followed by lengthy technical arguments justifying the interchanges of
summation signs. (In this aspect, our algebraic replica of the analysis of
infinite sums doesn't differ that much from the original.) We begin with the
computation; the justifying arguments will be sketched afterwards.

Let $f,g\in K\left[  \left[  x\right]  \right]  _{0}$. Thus, $\left[
x^{0}\right]  f=0$ and $\left[  x^{0}\right]  g=0$. Hence, $f+g\in K\left[
\left[  x\right]  \right]  _{0}$ (since (\ref{pf.thm.fps.ring.xn(a+b)=})
yields $\left[  x^{0}\right]  \left(  f+g\right)  =\underbrace{\left[
x^{0}\right]  f}_{=0}+\underbrace{\left[  x^{0}\right]  g}_{=0}=0$).

By the definition of $\operatorname*{Exp}$, we have%
\[
\operatorname*{Exp}f=\exp\circ f=\exp\left[  f\right]  =\sum_{n\in\mathbb{N}%
}\dfrac{1}{n!}f^{n}%
\]
(by Definition \ref{def.fps.subs}, since $\exp=\sum_{n\in\mathbb{N}}\dfrac
{1}{n!}x^{n}$). Similarly,%
\[
\operatorname*{Exp}g=\sum_{n\in\mathbb{N}}\dfrac{1}{n!}g^{n}%
\]
and%
\[
\operatorname*{Exp}\left(  f+g\right)  =\sum_{n\in\mathbb{N}}\dfrac{1}%
{n!}\left(  f+g\right)  ^{n}.
\]


Now, the latter equality becomes%
\begin{align*}
\operatorname*{Exp}\left(  f+g\right)   &  =\sum_{n\in\mathbb{N}}\dfrac{1}%
{n!}\underbrace{\left(  f+g\right)  ^{n}}_{\substack{=\sum\limits_{k=0}%
^{n}\dbinom{n}{k}f^{k}g^{n-k}\\\text{(by the binomial theorem)}}}=\sum
_{n\in\mathbb{N}}\dfrac{1}{n!}\sum\limits_{k=0}^{n}\dbinom{n}{k}f^{k}g^{n-k}\\
&  =\sum_{n\in\mathbb{N}}\ \ \sum\limits_{k=0}^{n}\underbrace{\dfrac{1}%
{n!}\dbinom{n}{k}}_{\substack{=\dfrac{1}{k!\left(  n-k\right)  !}\\\text{(by
(\ref{eq.binom.fac-form}))}}}f^{k}g^{n-k}=\underbrace{\sum_{n\in\mathbb{N}%
}\ \ \sum\limits_{k=0}^{n}}_{\substack{=\sum_{\substack{\left(  n,k\right)
\in\mathbb{N}\times\mathbb{N};\\k\leq n}}\\=\sum_{k\in\mathbb{N}}%
\ \ \sum\limits_{n\geq k}}}\dfrac{1}{k!\left(  n-k\right)  !}f^{k}g^{n-k}\\
&  =\sum_{k\in\mathbb{N}}\ \ \sum\limits_{n\geq k}\dfrac{1}{k!\left(
n-k\right)  !}f^{k}g^{n-k}=\sum_{k\in\mathbb{N}}\ \ \sum\limits_{\ell
\in\mathbb{N}}\dfrac{1}{k!\ell!}f^{k}g^{\ell}\\
&  \ \ \ \ \ \ \ \ \ \ \ \ \ \ \ \ \ \ \ \ \left(  \text{here, we have
substituted }\ell\text{ for }n-k\text{ in the second sum}\right)  .
\end{align*}
Comparing this with%
\begin{align*}
\operatorname*{Exp}f\cdot\operatorname*{Exp}g  &  =\left(  \sum_{k\in
\mathbb{N}}\dfrac{1}{k!}f^{k}\right)  \cdot\left(  \sum_{\ell\in\mathbb{N}%
}\dfrac{1}{\ell!}g^{\ell}\right) \\
&  \ \ \ \ \ \ \ \ \ \ \ \ \ \ \ \ \ \ \ \ \left(
\begin{array}
[c]{c}%
\text{since }\operatorname*{Exp}f=\sum_{n\in\mathbb{N}}\dfrac{1}{n!}f^{n}%
=\sum_{k\in\mathbb{N}}\dfrac{1}{k!}f^{k}\\
\text{and }\operatorname*{Exp}g=\sum_{n\in\mathbb{N}}\dfrac{1}{n!}g^{n}%
=\sum_{\ell\in\mathbb{N}}\dfrac{1}{\ell!}g^{\ell}%
\end{array}
\right) \\
&  =\sum_{k\in\mathbb{N}}\ \ \sum_{\ell\in\mathbb{N}}\dfrac{1}{k!}f^{k}%
\cdot\dfrac{1}{\ell!}g^{\ell}=\sum_{k\in\mathbb{N}}\ \ \sum_{\ell\in
\mathbb{N}}\dfrac{1}{k!\ell!}f^{k}g^{\ell},
\end{align*}
we obtain $\operatorname*{Exp}\left(  f+g\right)  =\operatorname*{Exp}%
f\cdot\operatorname*{Exp}g$.

\begin{fineprint}
This is sufficient to prove Lemma \ref{lem.fps.Exp-Log-additive} \textbf{(a)}
if we can justify the above manipulations of infinite sums. Actually, there is
just one manipulation that we need to justify, and that is our replacement of
\textquotedblleft$\sum_{n\in\mathbb{N}}\ \ \sum\limits_{k=0}^{n}%
$\textquotedblright\ by \textquotedblleft$\sum_{k\in\mathbb{N}}\ \ \sum
\limits_{n\geq k}$\textquotedblright. This is an application of the
\textquotedblleft discrete Fubini rule\textquotedblright\ (specifically, of a
version thereof in which the summation is over all pairs $\left(  n,k\right)
\in\mathbb{N}\times\mathbb{N}$ satisfying $k\leq n$). In order to justify this
manipulation, we need to show that the family $\left(  \dfrac{1}{k!\left(
n-k\right)  !}f^{k}g^{n-k}\right)  _{\left(  n,k\right)  \in\mathbb{N}%
\times\mathbb{N}\text{ satisfying }k\leq n}$ is summable. In other words, we
need to show the following statement:

\begin{statement}
\textit{Statement 1:} For each $m\in\mathbb{N}$, all but finitely many pairs
$\left(  n,k\right)  \in\mathbb{N}\times\mathbb{N}$ satisfying $k\leq n$
satisfy $\left[  x^{m}\right]  \left(  \dfrac{1}{k!\left(  n-k\right)  !}%
f^{k}g^{n-k}\right)  =0$.
\end{statement}

We shall achieve this by proving the following statement:

\begin{statement}
\textit{Statement 2:} For any three nonnegative integers $m,k,\ell$ with
$m<k+\ell$, we have $\left[  x^{m}\right]  \left(  f^{k}g^{\ell}\right)  =0$.
\end{statement}

[\textit{Proof of Statement 2:} Let $m,k,\ell$ be three nonnegative integers
with $m<k+\ell$. We must show that $\left[  x^{m}\right]  \left(  f^{k}%
g^{\ell}\right)  =0$.

We have $\left[  x^{0}\right]  f=0$. Hence, Lemma \ref{lem.fps.g=xh} (applied
to $a=f$) shows that there exists an $h\in K\left[  \left[  x\right]  \right]
$ such that $f=xh$. Consider this $h$ and denote it by $u$. Thus, $u\in
K\left[  \left[  x\right]  \right]  $ and $f=xu$.

We have $\left[  x^{0}\right]  g=0$. Hence, Lemma \ref{lem.fps.g=xh} (applied
to $a=g$) shows that there exists an $h\in K\left[  \left[  x\right]  \right]
$ such that $g=xh$. Consider this $h$ and denote it by $v$. Thus, $v\in
K\left[  \left[  x\right]  \right]  $ and $g=xv$.

Now, from $f=xu$ and $g=xv$, we obtain $f^{k}g^{\ell}=\left(  xu\right)
^{k}\left(  xv\right)  ^{\ell}=x^{k}u^{k}x^{\ell}v^{\ell}=x^{k+\ell}%
u^{k}v^{\ell}$. However, Lemma \ref{lem.fps.first-n-coeffs-of-xna} (applied to
$k+\ell$ and $u^{k}v^{\ell}$ instead of $k$ and $a$) shows that the first
$k+\ell$ coefficients of the FPS $x^{k+\ell}u^{k}v^{\ell}$ are $0$. In other
words, the first $k+\ell$ coefficients of the FPS $f^{k}g^{\ell}$ are $0$
(since $f^{k}g^{\ell}=x^{k+\ell}u^{k}v^{\ell}$). But $\left[  x^{m}\right]
\left(  f^{k}g^{\ell}\right)  $ is one of these first $k+\ell$ coefficients
(since $m<k+\ell$). Thus, we conclude that $\left[  x^{m}\right]  \left(
f^{k}g^{\ell}\right)  =0$. This proves Statement 2.]

Note that Statement 2 entails that the family $\left(  f^{k}g^{\ell}\right)
_{\left(  k,\ell\right)  \in\mathbb{N}\times\mathbb{N}}$ is summable (because
when $m\in\mathbb{N}$ is given, all but finitely many pairs $\left(
k,\ell\right)  \in\mathbb{N}\times\mathbb{N}$ satisfy $m<k+\ell$). However, we
need to prove Statement 1, so let us do this:

[\textit{Proof of Statement 1:} Let $m\in\mathbb{N}$. If $\left(  n,k\right)
\in\mathbb{N}\times\mathbb{N}$ is a pair satisfying $k\leq n$ and $m<n$, then
\begin{align*}
\left[  x^{m}\right]  \left(  \dfrac{1}{k!\left(  n-k\right)  !}f^{k}%
g^{n-k}\right)   &  =\dfrac{1}{k!\left(  n-k\right)  !}\underbrace{\left[
x^{m}\right]  \left(  f^{k}g^{n-k}\right)  }_{\substack{=0\\\text{(by
Statement 2}\\\text{(applied to }\ell=n-k\text{),}\\\text{since }m<n=k+\left(
n-k\right)  \text{)}}}\ \ \ \ \ \ \ \ \ \ \left(  \text{by
(\ref{pf.thm.fps.ring.xn(la)=})}\right) \\
&  =0.
\end{align*}
Thus, all but finitely many pairs $\left(  n,k\right)  \in\mathbb{N}%
\times\mathbb{N}$ satisfying $k\leq n$ satisfy \newline$\left[  x^{m}\right]
\left(  \dfrac{1}{k!\left(  n-k\right)  !}f^{k}g^{n-k}\right)  =0$ (because
all but finitely many such pairs satisfy $m<n$). This proves Statement 1.]

As explained above, Statement 1 shows that the family \newline$\left(
\dfrac{1}{k!\left(  n-k\right)  !}f^{k}g^{n-k}\right)  _{\left(  n,k\right)
\in\mathbb{N}\times\mathbb{N}\text{ satisfying }k\leq n}$ is summable, and
thus our interchange of summation signs made above is justified. This
completes our proof of Lemma \ref{lem.fps.Exp-Log-additive} \textbf{(a)}.
\end{fineprint}

\bigskip

\textbf{(b)} This easily follows from part \textbf{(a)}, since we know that
$\operatorname*{Log}$ is inverse to $\operatorname*{Exp}$. Here are the details:

Let $f,g\in K\left[  \left[  x\right]  \right]  _{1}$. Set
$u=\operatorname*{Log}f$ and $v=\operatorname*{Log}g$; then, $u,v\in K\left[
\left[  x\right]  \right]  _{0}$ (since $\operatorname*{Log}$ is a map from
$K\left[  \left[  x\right]  \right]  _{1}$ to $K\left[  \left[  x\right]
\right]  _{0}$). Hence, Lemma \ref{lem.fps.Exp-Log-additive} \textbf{(a)}
(applied to $u$ and $v$ instead of $f$ and $g$) yields $\operatorname*{Exp}%
\left(  u+v\right)  =\operatorname*{Exp}u\cdot\operatorname*{Exp}v$.

However, Lemma \ref{lem.fps.Exp-Log-maps-inv} says that the maps
$\operatorname*{Exp}$ and $\operatorname*{Log}$ are mutually inverse
bijections between $K\left[  \left[  x\right]  \right]  _{0}$ and $K\left[
\left[  x\right]  \right]  _{1}$. Hence, $\operatorname*{Exp}\circ
\operatorname*{Log}=\operatorname*{id}$ and $\operatorname*{Log}%
\circ\operatorname*{Exp}=\operatorname*{id}$.

Now, from $u=\operatorname*{Log}f$, we obtain $\operatorname*{Exp}%
u=\operatorname*{Exp}\left(  \operatorname*{Log}f\right)  =\underbrace{\left(
\operatorname*{Exp}\circ\operatorname*{Log}\right)  }_{=\operatorname*{id}%
}\left(  f\right)  =\operatorname*{id}\left(  f\right)  =f$. Similarly,
$\operatorname*{Exp}v=g$. Multiplying these two equalities, we find
$\operatorname*{Exp}u\cdot\operatorname*{Exp}v=fg$. Now, we have%
\[
\operatorname*{Log}\left(  \operatorname*{Exp}\left(  u+v\right)  \right)
=\underbrace{\left(  \operatorname*{Log}\circ\operatorname*{Exp}\right)
}_{=\operatorname*{id}}\left(  u+v\right)  =\operatorname*{id}\left(
u+v\right)  =u+v=\operatorname*{Log}f+\operatorname*{Log}g
\]
(since $u=\operatorname*{Log}f$ and $v=\operatorname*{Log}g$). In view of
$\operatorname*{Exp}\left(  u+v\right)  =\operatorname*{Exp}u\cdot
\operatorname*{Exp}v=fg$, this rewrites as $\operatorname*{Log}\left(
fg\right)  =\operatorname*{Log}f+\operatorname*{Log}g$. This proves Lemma
\ref{lem.fps.Exp-Log-additive} \textbf{(b)}.
\end{proof}

We can neatly pack the last few lemmas into a single theorem through the use
of group isomorphisms. To this purpose, we need to observe that $K\left[
\left[  x\right]  \right]  _{0}$ is a group under addition and $K\left[
\left[  x\right]  \right]  _{1}$ is a group under multiplication:

\begin{proposition}
\label{prop.fps.Exp-Log-groups}\textbf{(a)} The subset $K\left[  \left[
x\right]  \right]  _{0}$ of $K\left[  \left[  x\right]  \right]  $ is closed
under addition and subtraction and contains $0$, and thus forms a group
$\left(  K\left[  \left[  x\right]  \right]  _{0},+,0\right)  $. \medskip

\textbf{(b)} The subset $K\left[  \left[  x\right]  \right]  _{1}$ of
$K\left[  \left[  x\right]  \right]  $ is closed under multiplication and
division and contains $1$, and thus forms a group $\left(  K\left[  \left[
x\right]  \right]  _{1},\cdot,1\right)  $.
\end{proposition}

\begin{proof}
[Proof of Proposition \ref{prop.fps.Exp-Log-groups}.]\textbf{(a)} It is clear
that the set $K\left[  \left[  x\right]  \right]  _{0}$ contains the FPS $0$
(since $\left[  x^{0}\right]  0=0$). Thus, it remains to show that $K\left[
\left[  x\right]  \right]  _{0}$ is closed under addition and subtraction. But
this is easy: If $f,g\in K\left[  \left[  x\right]  \right]  _{0}$, then
$\left[  x^{0}\right]  f=0$ and $\left[  x^{0}\right]  g=0$, and therefore
$f+g\in K\left[  \left[  x\right]  \right]  _{0}$ (since
(\ref{pf.thm.fps.ring.xn(a+b)=}) yields $\left[  x^{0}\right]  \left(
f+g\right)  =\underbrace{\left[  x^{0}\right]  f}_{=0}+\underbrace{\left[
x^{0}\right]  g}_{=0}=0$) and $f-g\in K\left[  \left[  x\right]  \right]
_{0}$ (by a similar argument using (\ref{pf.thm.fps.ring.xn(a-b)=})). Thus,
$K\left[  \left[  x\right]  \right]  _{0}$ is closed under addition and
subtraction. This proves Proposition \ref{prop.fps.Exp-Log-groups}
\textbf{(a)}.

\textbf{(b)} Any $a\in K\left[  \left[  x\right]  \right]  _{1}$ is invertible
in $K\left[  \left[  x\right]  \right]  $ (indeed, $a\in K\left[  \left[
x\right]  \right]  _{1}$ shows that $\left[  x^{0}\right]  a=1$; thus,
$\left[  x^{0}\right]  a$ is invertible in $K$; therefore, Proposition
\ref{prop.fps.invertible} entails that $a$ is invertible in $K\left[  \left[
x\right]  \right]  $). Hence, $\dfrac{f}{g}$ is well-defined for any $f,g\in
K\left[  \left[  x\right]  \right]  _{1}$.

Next, we claim that $K\left[  \left[  x\right]  \right]  _{1}$ is closed under
multiplication. Indeed, if $f,g\in K\left[  \left[  x\right]  \right]  _{1}$,
then $\left[  x^{0}\right]  f=1$ and $\left[  x^{0}\right]  g=1$, and
therefore $fg\in K\left[  \left[  x\right]  \right]  _{1}$ (since
(\ref{pf.thm.fps.ring.x0(ab)=}) yields $\left[  x^{0}\right]  \left(
fg\right)  =\underbrace{\left[  x^{0}\right]  f}_{=1}\cdot\underbrace{\left[
x^{0}\right]  g}_{=1}=1$). This shows that $K\left[  \left[  x\right]
\right]  _{1}$ is closed under multiplication.

It remains to prove that $K\left[  \left[  x\right]  \right]  _{1}$ is closed
under division. Indeed, if $f,g\in K\left[  \left[  x\right]  \right]  _{1}$,
then $\left[  x^{0}\right]  f=1$ and $\left[  x^{0}\right]  g=1$, and
therefore $\dfrac{f}{g}\in K\left[  \left[  x\right]  \right]  _{1}$ (because
we have $f=\dfrac{f}{g}\cdot g$ and thus
\begin{align*}
\left[  x^{0}\right]  f  &  =\left[  x^{0}\right]  \left(  \dfrac{f}{g}\cdot
g\right)  =\left[  x^{0}\right]  \dfrac{f}{g}\cdot\underbrace{\left[
x^{0}\right]  g}_{=1}\ \ \ \ \ \ \ \ \ \ \left(  \text{by
(\ref{pf.thm.fps.ring.x0(ab)=})}\right) \\
&  =\left[  x^{0}\right]  \dfrac{f}{g}%
\end{align*}
and thus $\left[  x^{0}\right]  \dfrac{f}{g}=\left[  x^{0}\right]  f=1$, so
that $\dfrac{f}{g}\in K\left[  \left[  x\right]  \right]  _{1}$). This shows
that $K\left[  \left[  x\right]  \right]  _{1}$ is closed under division.
Thus, Proposition \ref{prop.fps.Exp-Log-groups} \textbf{(b)} is proven.
\end{proof}

The two groups in Proposition \ref{prop.fps.Exp-Log-groups} can now be
connected through $\operatorname*{Exp}$ and $\operatorname*{Log}$:

\begin{theorem}
\label{thm.fps.Exp-Log-group-iso}The maps
\[
\operatorname*{Exp}:\left(  K\left[  \left[  x\right]  \right]  _{0}%
,+,0\right)  \rightarrow\left(  K\left[  \left[  x\right]  \right]  _{1}%
,\cdot,1\right)
\]
and%
\[
\operatorname*{Log}:\left(  K\left[  \left[  x\right]  \right]  _{1}%
,\cdot,1\right)  \rightarrow\left(  K\left[  \left[  x\right]  \right]
_{0},+,0\right)
\]
are mutually inverse group isomorphisms.
\end{theorem}

\begin{proof}
[Proof of Theorem \ref{thm.fps.Exp-Log-group-iso} (sketched).]Lemma
\ref{lem.fps.Exp-Log-additive} yields that these two maps are group
homomorphisms\footnote{Here, we are using the following fact: If $\left(
G,\ast,e_{G}\right)  $ and $\left(  H,\ast,e_{H}\right)  $ are any two groups,
and if $\Phi:G\rightarrow H$ is a map such that every $f,g\in G$ satisfy
$\Phi\left(  f\ast g\right)  =\Phi\left(  f\right)  \ast\Phi\left(  g\right)
$, then $\Phi$ is a group homomorphism.}. Lemma \ref{lem.fps.Exp-Log-maps-inv}
shows that they are mutually inverse. Combining these results, we conclude
that these two maps are mutually inverse group isomorphisms. This proves
Theorem \ref{thm.fps.Exp-Log-group-iso}.
\end{proof}

Theorem \ref{thm.fps.Exp-Log-group-iso} helps us turn addition into
multiplication and vice versa when it comes to FPSs, at least if the constant
terms are the right ones. This will come useful rather soon.

\subsubsection{The logarithmic derivative}

For future use, we shall define and briefly study one more concept related to logarithms:

\begin{definition}
\label{def.fps.loder.1}In this definition, we do \textbf{not} use Convention
\ref{conv.fps.exp.K-Q-alg}; thus, $K$ can be an arbitrary commutative ring.
However, we set $K\left[  \left[  x\right]  \right]  _{1}=\left\{  f\in
K\left[  \left[  x\right]  \right]  \ \mid\ \left[  x^{0}\right]  f=1\right\}
$.

For any FPS $f\in K\left[  \left[  x\right]  \right]  _{1}$, we define the
\emph{logarithmic derivative} $\operatorname*{loder}f\in K\left[  \left[
x\right]  \right]  $ to be the FPS $\dfrac{f^{\prime}}{f}$. (This is
well-defined, since $f$ is easily seen to be invertible.\footnotemark)
\end{definition}

\footnotetext{Indeed: Let $f\in K\left[  \left[  x\right]  \right]  _{1}$.
Then, $\left[  x^{0}\right]  f=1$. Thus, $\left[  x^{0}\right]  f$ is
invertible. Hence, Proposition \ref{prop.fps.invertible} (applied to $a=f$)
yields that $f$ is invertible.}The reason why this FPS $\operatorname*{loder}%
f$ is called the logarithmic derivative of $f$ is made clear by the following
simple fact:

\begin{proposition}
\label{prop.fps.loder.log}Let $K$ be a commutative $\mathbb{Q}$-algebra. Let
$f\in K\left[  \left[  x\right]  \right]  _{1}$ be an FPS. Then,
$\operatorname*{loder}f=\left(  \operatorname*{Log}f\right)  ^{\prime}$.
\end{proposition}

\begin{proof}
[Proof of Proposition \ref{prop.fps.loder.log}.]The definition of
$\operatorname*{Log}$ yields $\operatorname*{Log}f=\overline{\log}\circ\left(
f-1\right)  $.

From $f\in K\left[  \left[  x\right]  \right]  _{1}$, we obtain $\left[
x^{0}\right]  f=1=\left[  x^{0}\right]  1$. Now, $\left[  x^{0}\right]
\left(  f-1\right)  =\left[  x^{0}\right]  f-\left[  x^{0}\right]  1=0$ (since
$\left[  x^{0}\right]  f=\left[  x^{0}\right]  1$). Hence, Proposition
\ref{prop.fps.exp-log-der} \textbf{(b)} (applied to $g=f-1$) yields
\[
\left(  \overline{\log}\circ\left(  f-1\right)  \right)  ^{\prime}=\left(
\underbrace{1+\left(  f-1\right)  }_{=f}\right)  ^{-1}\cdot\left(  f-1\right)
^{\prime}=f^{-1}\cdot\left(  f-1\right)  ^{\prime}=\dfrac{\left(  f-1\right)
^{\prime}}{f}.
\]
In view of $\operatorname*{Log}f=\overline{\log}\circ\left(  f-1\right)  $, we
can rewrite this as%
\[
\left(  \operatorname*{Log}f\right)  ^{\prime}=\dfrac{\left(  f-1\right)
^{\prime}}{f}.
\]


Since $\left(  f-1\right)  ^{\prime}=f^{\prime}$ (which is easy to
see\footnote{\textit{Proof.} Theorem \ref{thm.fps.deriv.rules} \textbf{(a)}
(applied to $g=-1$) yields $\left(  f+\left(  -1\right)  \right)  ^{\prime
}=f^{\prime}+\underbrace{\left(  -1\right)  ^{\prime}}_{=0}=f^{\prime}$. In
other words, $\left(  f-1\right)  ^{\prime}=f^{\prime}$.}), we can rewrite
this further as $\left(  \operatorname*{Log}f\right)  ^{\prime}=\dfrac
{f^{\prime}}{f}=\operatorname*{loder}f$ (since $\operatorname*{loder}f$ is
defined as $\dfrac{f^{\prime}}{f}$). This proves Proposition
\ref{prop.fps.loder.log}.
\end{proof}

Note that Proposition \ref{prop.fps.loder.log} only makes sense when $K$ is a
$\mathbb{Q}$-algebra (otherwise, $\operatorname*{Log}f$ would make no sense),
which is why we are not using it as a definition of $\operatorname*{loder}f$.
The logarithmic derivative is defined even when the logarithm is not!

If you have seen the logarithmic derivative in analysis, you will likely
expect the following property:

\begin{proposition}
\label{prop.fps.loder.prod}Let $f,g\in K\left[  \left[  x\right]  \right]
_{1}$ be two FPSs. Then, $\operatorname*{loder}\left(  fg\right)
=\operatorname*{loder}f+\operatorname*{loder}g$.

(Here, we do \textbf{not} use Convention \ref{conv.fps.exp.K-Q-alg}; thus, $K$
can be an arbitrary commutative ring.)
\end{proposition}

\begin{proof}
[Proof of Proposition \ref{prop.fps.loder.prod}.]The definition of a
logarithmic derivative yields $\operatorname*{loder}f=\dfrac{f^{\prime}}{f}$
and $\operatorname*{loder}g=\dfrac{g^{\prime}}{g}$, but it also yields%
\begin{align*}
\operatorname*{loder}\left(  fg\right)   &  =\dfrac{\left(  fg\right)
^{\prime}}{fg}=\dfrac{f^{\prime}g+fg^{\prime}}{fg}\ \ \ \ \ \ \ \ \ \ \left(
\begin{array}
[c]{c}%
\text{since Theorem \ref{thm.fps.deriv.rules} \textbf{(d)}}\\
\text{says that }\left(  fg\right)  ^{\prime}=f^{\prime}g+fg^{\prime}%
\end{array}
\right) \\
&  =\underbrace{\dfrac{f^{\prime}}{f}}_{=\operatorname*{loder}f}%
+\underbrace{\dfrac{g^{\prime}}{g}}_{=\operatorname*{loder}g}%
=\operatorname*{loder}f+\operatorname*{loder}g.
\end{align*}
This proves Proposition \ref{prop.fps.loder.prod}.
\end{proof}

\begin{corollary}
\label{cor.fps.loder.prodk}Let $f_{1},f_{2},\ldots,f_{k}$ be any $k$ FPSs in
$K\left[  \left[  x\right]  \right]  _{1}$. Then,
\[
\operatorname*{loder}\left(  f_{1}f_{2}\cdots f_{k}\right)
=\operatorname*{loder}\left(  f_{1}\right)  +\operatorname*{loder}\left(
f_{2}\right)  +\cdots+\operatorname*{loder}\left(  f_{k}\right)  .
\]


(Here, we do \textbf{not} use Convention \ref{conv.fps.exp.K-Q-alg}; thus, $K$
can be an arbitrary commutative ring.)
\end{corollary}

\begin{proof}
[Proof of Corollary \ref{cor.fps.loder.prodk}.]Induct on $k$. The \textit{base
case} ($k=0$) requires showing that $\operatorname*{loder}1=0$, which is easy
to see from the definition of a logarithmic derivative (since $1^{\prime}=0$).
The \textit{induction step} follows easily from Proposition
\ref{prop.fps.loder.prod}.
\end{proof}

\begin{corollary}
\label{cor.fps.loder.inv}Let $f$ be any FPS in $K\left[  \left[  x\right]
\right]  _{1}$. Then, $\operatorname*{loder}\left(  f^{-1}\right)
=-\operatorname*{loder}f$.

(Here, we do \textbf{not} use Convention \ref{conv.fps.exp.K-Q-alg}; thus, $K$
can be an arbitrary commutative ring.)
\end{corollary}

\begin{proof}
[Proof of Corollary \ref{cor.fps.loder.inv}.]First of all, $f$ is invertible
(as we showed in Definition \ref{def.fps.loder.1}). This shows that $f^{-1}$
is well-defined.

Proposition \ref{prop.fps.loder.prod} (applied to $g=f^{-1}$) yields
$\operatorname*{loder}\left(  ff^{-1}\right)  =\operatorname*{loder}%
f+\operatorname*{loder}\left(  f^{-1}\right)  $. However, we also have
$\operatorname*{loder}\left(  \underbrace{ff^{-1}}_{=1}\right)
=\operatorname*{loder}1=0$ (since $1^{\prime}=0$). Comparing these two
equalities, we obtain $\operatorname*{loder}f+\operatorname*{loder}\left(
f^{-1}\right)  =0$. In other words, $\operatorname*{loder}\left(
f^{-1}\right)  =-\operatorname*{loder}f$. This proves Corollary
\ref{cor.fps.loder.inv}.
\end{proof}
